% !PRINTED 373
\weberpage{373}
%<<CONT
% Source-faithful Cayley configuration signs: three/four concurrent strokes and a seven-rayed star.
\providecommand{\weberthreeup}{\mathord{\tikz[baseline=-0.2ex,x=0.23em,y=0.23em,line width=0.45pt]{\draw (0,0)--(-1,1) (0,0)--(0,1.08) (0,0)--(1,1);}}}
\providecommand{\weberfourup}{\mathord{\tikz[baseline=-0.2ex,x=0.23em,y=0.23em,line width=0.45pt]{\draw (0,0)--(-1.05,1) (0,0)--(-0.35,1.08) (0,0)--(0.35,1.08) (0,0)--(1.05,1);}}}
\providecommand{\webersevenstar}{\mathord{\tikz[baseline=-0.35ex,x=0.21em,y=0.21em,line width=0.45pt]{\draw (0,0)--(90:1) (0,0)--(38.57:1) (0,0)--(-12.86:1) (0,0)--(-64.29:1) (0,0)--(-115.71:1) (0,0)--(-167.14:1) (0,0)--(141.43:1);}}}
\gesperrt{Die Quadrupel von Doppeltangenten mit den
Zeichen}
\[
\mid\mid\mid\mid,\ \square
\]
\gesperrt{sind syzygetisch, und die mit den Zeichen}
\[
\weberfourup,\ \triangle\mid,\ \weberthreeup\mid,\ \vee\vee
\]
\gesperrt{azygetisch.}

Alle übrigen Quadrupel sind weder syzygetisch noch azygetisch
.

Hiernach ist es leicht, die Zeichen für sämmtliche vollständige
 Siebener-Systeme zu bilden.

Ein solches Zeichen muss aus sieben Strichen bestehen, die
nicht mehr als acht Eckpunkte haben können und die eine Figur
bilden, aus der sich keine der beiden Figuren $\mid\mid\mid$, $\sqcap$ ablösen
lässt. Daraus folgt zunächst, dass diese Figur aus nicht mehr als
zwei getrennten Theilen bestehen kann, weil sonst die Figur $\mid\mid\mid$
darin enthalten wäre, und dass kein Theil mehr als ein Centrum
haben kann, von dem mehrere Striche auslaufen, ausser wenn
dieser Theil das Dreieck $\triangle$ ist, weil sonst die Figur $\sqcap$ vorkommen
 würde.

Wenn nun die Figur eintheilig ist, so muss sie ein siebenstrahliger
 Stern $\webersevenstar$ sein, und da man jede der acht Ziffern als
Mittelpunkt wählen kann, so sind dies acht Möglichkeiten, von
denen eine die oben betrachtete Annahme ist:
\[
[1\,8]\,[2\,8]\,[3\,8]\,[4\,8]\,[5\,8]\,[6\,8]\,[7\,8].
\]
Ist aber die Figur zweitheilig, so ist zunächst auszuschliessen,
dass der eine Theil aus einem oder aus zwei Strichen oder
einem dreistrahligen Sterne besteht; denn in diesen Fällen müsste
der andere Theil ein Stern mit sechs, fünf oder vier Strahlen
sein. Dazu aber bleiben von den acht Ziffern nicht mehr genug
übrig. Es bleibt also nur noch übrig, dass der eine Theil ein
Dreieck, der andere ein vierstrahliger Stern ist, $\triangle\,\weberfourup$, und diese
Annahme ist auch in der That immer zulässig. Ein Repräsentant
eines solchen Systems ist:
\[
[1\,2]\,[1\,3]\,[2\,3]\,[4\,5]\,[4\,6]\,[4\,7]\,[4\,8].
\]
Das Dreieck kann man auf $8.7.6 : 2.3 = 56$ verschiedene
Arten wählen, und zu jedem Dreieck giebt es noch fünf Möglichkeiten
, den vierstrahligen Stern anzunehmen, da man jeden der
übrigen fünf Punkte zum Centrum machen kann. Die Anzahl
%CONT>>
